\section{The rational newform attached to the type-II corner}

% BEGIN R72Q1
\paragraph{R72Q1. The given corner stays irreducible over the marked field.}
Use E72G1 to fix the AFE15 link, its mark $g_1$ and its seed. Take an
original finite coefficient DVR, ordinary point $r$, and type-II
extension $W:G_{\mathbb Q}\to{\rm GL}_2(\overline{\mathbb Q}_\ell)$
as in $\mathrm{AutEnd2}$. Write $H=\ker\eta=G_K$ and $J=G_L$.
Then $J\subset H$, $L/\mathbb Q$ is abelian and the marked element
lies in $J$ with $\bar\Psi(g_1)^2\ne1$. Indeed its image is the
frame's distinguished element, and the frame's base-subgroup
condition gives precisely this inequality. The extension $W$ is
continuous: on the open subgroup $H$ its four matrix entries are
entries of the continuous ordinary point in the given finite DVR,
and continuity on all cosets follows by translation. Type II
gives irreducibility on $H$.

The corner has residual trace $\bar\Psi+\bar\Psi^{-1}$ on $H$.
If its restriction to $J$ were reducible, there would be a
nontrivial character $\nu$ of $H/J$ with
$W|_H\simeq W|_H\otimes\nu$. Indeed the restriction to the finite
index normal subgroup $J$ is semisimple by Clifford theory
(the sum of the translates of an irreducible subspace is the
whole irreducible module). Its commuting algebra would then have
dimension greater than one. The finite abelian group $H/J$ acts
on this algebra by conjugation; its invariant subspace consists
only of the scalars. A nontrivial character line therefore gives
an intertwiner, necessarily an isomorphism by irreducibility.
Choose $g\in H$ with $\nu(g)\ne1$. Both $g$ and $gg_1$ have
zero trace. Reducing the two trace identities gives
$\bar\Psi(g)^2=\bar\Psi(gg_1)^2=-1$, contrary to the marked
inequality. Thus $W|_{G_L}$ is absolutely irreducible.

The type-II zero entries make the coordinate plane with indices
$0,2$ an invariant block on $H$. Consequently its corner is a
homomorphism there over the original DVR. Restrict the ordinary
point to the linked quotient and take this block. More explicitly
the ordinary-point construction on the frame gives its three
dimensional representation, whose matrices on the even subgroup
are block matrices by surjectivity of the link projection and
the type-II equations. Taking the common $0,2$ block gives the
source's rank-two lift in $\mathrm{exists\_nrep2}$. Its residual
diagonal characters are $\bar\Psi,\bar\Psi^{-1}$ and its inertia
characters are exactly the corresponding $\Psi^{w(2,i)}$, as
required by $\mathrm{AutPtAlt}$. The matrix identity for the
prescribed $W$ gives equality on every element of $G_L$ after
including coefficients. The map of finite DVRs from the original
coefficient ring is local by integrality. Invoke the rank-two
output of L69D3 on this lift and the irreducibility just proved.
It gives a regular algebraic conjugate-self-dual cusp over $L$
attached, on geometric Frobenius in the $\mathrm{AttC}$
normalization, to $W|_{G_L}\otimes a_2$, with the same algebraic
character $a_2$ furnished by that link.
% END R72Q1

% BEGIN R72Q2
\paragraph{R72Q2. The two jumps and descent to a rational cusp.}
Write $\epsilon$ for the $\ell$-adic cyclotomic character. We
compute the jumps of $W$ using the open part of inertia
$I_0=I_\ell\cap\ker\eta\cap\ker(\mathrm{dirG}(\psi))$ at the
specified valuation. Here $\Psi=\epsilon^e$. On this subgroup
$r$ is ordinary with three inertia characters
$\epsilon^e,1,\epsilon^{-e}$, and its middle coordinate line
is trivial by the stipulated $(1,1)$ entry. Project its invariant
first ordinary line to the $0,2$ plane. The projection is nonzero
since $\epsilon^e$ is nontrivial on $I_0$ and the middle line is
trivial. Indeed $I_0$ has finite index in inertia, the cyclotomic
image of rational $\ell$-adic inertia is infinite and $e>0$.
The projection thus gives an $\epsilon^e$ line for the prescribed
$W$ on $I_0$. The other character is $\epsilon^{-e}$ by the
three ordinary diagonal characters and the middle coordinate.
These two characters are distinct on $I_0$; the
$\epsilon^e$-eigenspace is exactly a line. The group $I_0$ is
normal in the decomposition group, and both characters extend
to the latter. Hence its $\epsilon^e$-line is decomposition-group
stable. On an open local subgroup whose inertia is contained
in $I_0$, the two factors differ from the specified cyclotomic
powers only by unramified characters.

The rank-two automorphic representation in R72Q1 has a de Rham
attachment by A72F8: a conjugate-self-dual cusp in this even
rank admits the trivial polarizing character. The attachment is
the given twist by Chebotarev and irreducibility. Its algebraic
character factor is de Rham by A72F3. Consequently $W$ is de
Rham after restriction at places over $\ell$ in $L$. The
de Rham property descends under finite extension by Galois
descent for $D_{\rm dR}$, after taking a local normal closure.
The de Rham functor is exact and strict on de Rham representations,
and the two rank-one factors on the open subgroup just constructed
have jumps $-e,e$. Thus $V=W\epsilon^{-e}$ has jumps $\{0,2e\}$
and remains absolutely irreducible on $G_L$.

Untwist the cusp over $L$ by the Hecke character belonging to
$a_2^{-1}$ and twist by the algebraic Hecke character belonging
to $\epsilon^{-e}$. Rank-one reciprocity and multiplication of
Satake eigenvalues give a regular algebraic cusp attached to
$V|_{G_L}$. At each complex embedding its infinitesimal exponents
are $\{1/2,1/2-2e\}$. Indeed A72F8 identifies the exponents for
the original conjugate-self-dual cusp with $1/2$ minus its
labelled jumps, and A72F3 gives the negative infinity exponent
as jump for each character twist. Jumps of a restriction of a
de Rham representation over $\mathbb Q_\ell$ are unchanged at
every embedding of the extension.

Descend first from $L$ to the imaginary quadratic $K$ through
a series with cyclic prime-degree steps. Each intermediate field
is CM because $L/\mathbb Q$ is abelian and the intermediate field
contains $K$. D70W1 applies at every step with the same $V$,
irreducible already on $G_L$. We obtain a cusp $\Pi_K$ attached
to $V|_{G_K}$ with the same infinitesimal exponents. It is fixed
by $\operatorname{Gal}(K/\mathbb Q)$: conjugate Frobenius elements
give equal multisets and strong multiplicity one applies. The
cyclic descent theorem of Arthur--Clozel (\emph{Simple algebras,
base change, and the advanced theory of the trace formula}, Ch.~III) gives a cuspidal representation $\Pi$ over
$\mathbb Q$ with base change $\Pi_K$. At the archimedean place
the restriction of its Weil parameter to $\mathbb C^\times$ is
that of $\Pi_K$.
% END R72Q2

% BEGIN R72Q3
\paragraph{R72Q3. Weight, conductor and the exact twist of the given extension.}
For the rational cusp in R72Q2 the two angular possibilities
can be distinguished at infinity. If its real Weil parameter
is reducible, the two characters restrict to characters of
$\mathbb C^\times$ fixed by complex conjugation, so their
two infinitesimal exponents are paired with themselves. The
normalized inducing characters of $\mathbb R^\times$ then
have real exponents differing by $2e\ge2$. A cuspidal
$\mathrm{GL}_2$ representation is globally generic and
unitarizable after a common norm twist. The unitary dual of
$\mathrm{GL}_2(\mathbb R)$ allows a generic irreducible principal
series with unequal real exponents only in the complementary
range where their difference is less than one; a finite-dimensional
quotient at a reducibility point is nongeneric. This excludes the
reducible parameter. Thus it is induced from the non-conjugation-fixed
character of $\mathbb C^\times$. Its angular exponent has absolute
value $2e$, and the discrete series has lowest holomorphic weight
$k=2e+1$ (the real local Langlands classification).
The common real shift of the two infinitesimal exponents is
$1/2-e$. Twisting $\Pi$ by $|\det|^{e-1/2}$ yields the unitary
holomorphic discrete-series normalization.

Take the normalized holomorphic newvector at its conductor $M$
using the classical newform correspondence and the local
newvector theorem. This gives a cusp form on $\Gamma_1(M)$
with nebentypus, first Fourier coefficient one, and the usual
Hecke coefficient recursion at every prime away from $M$.
If a nonzero cusp form at a proper divisor of $M$ has the same
cofinite Hecke eigenvalues, project it to a nonzero cuspidal
constituent and take the primitive eigenform of that constituent.
Its nebentypus may differ initially. The specified eigenvalues
give equality of Frobenius traces of the two semisimple Deligne
representations at all but finitely many primes. Chebotarev and
Brauer--Nesbitt identify the representations and hence also their
determinants. The full good-prime Hecke polynomials now agree, so
strong multiplicity one identifies the two cusps. This contradicts
minimality of the conductor. Thus it meets the precise
$\mathrm{Newform.new}$ test, including old vectors.

Let $s_f$ be the semisimple Deligne representation of this form
in the geometric Frobenius convention (take the dual of the
arithmetic-Frobenius convention). Its geometric Frobenius roots
at a good prime $p$ are those of the unitary cusp multiplied
by $p^{(k-1)/2}$. The common norm twist above multiplied
the roots of $\Pi_p$ by $p^{-(e-1/2)}$, so the net factor is
$p^{1/2}$, exactly $\mathrm{AttC}$ for $\Pi$. On restriction
to $K$ both the Satake roots and this factor are raised to
the local residue degree. Hence $s_f|_{G_K}\simeq V|_{G_K}$
by Chebotarev, Brauer--Nesbitt and the irreducibility on $G_K$.
Both representations are defined on the full rational group.
After identifying the restrictions, Schur's lemma makes their
difference a character of $\operatorname{Gal}(K/\mathbb Q)$.
Twist the rational cusp by that quadratic character if needed
and take its newvector at its own conductor. The weight is
unchanged and the geometric Galois representation is now the
specified $V$ up to isomorphism. Thus, at cofinitely many primes
and for every prolongation, inertia is trivial and the polynomial
of inverse arithmetic Frobenius is
\[
 X^2-a_p(f)X+\chi_f(p)p^{\,2e}.
\]
Semisimplicity follows from irreducibility. This is
$\mathrm{Newform.Att}$ for $W\epsilon^{-e}$ at the original
$\iota$. Every step applies to every allowed type-II $r$ and
every extension $W$ with the stated middle-coordinate and
trace-zero conditions at the already fixed $g_1$. No level
is prescribed by $\mathrm{AutEnd2}$. This supplies that endpoint
from the rank-two output of L69D3 on the AFE15 frame.
% END R72Q3
